\section{Arquitectura del Modelo}

El modelo \texttt{SentimentRNN} es una red neuronal recurrente apilada configurable,
diseñada para clasificación binaria de sentimiento. Su arquitectura soporta tanto
celdas LSTM como GRU como tipo de capa recurrente, permitiendo una comparación
directa entre ambas variantes bajo condiciones idénticas. A continuación se describe
cada componente de la arquitectura, cuyo resumen se presenta en la
Tabla~\ref{tab:arquitectura}.

\subsection{Capa de Embedding}

La primera capa del modelo es un embedding que transforma cada token del vocabulario
en un vector denso de representación distribuida. El vocabulario tiene un tamaño de
25,000 tokens y cada uno se proyecta a un espacio de dimensión 128. Esta capa
contiene $25{,}000 \times 128 = 3{,}200{,}000$ parámetros entrenables y constituye
la mayor parte de los parámetros del modelo.

\subsection{Capas Recurrentes Apiladas}

El modelo utiliza 2 capas recurrentes apiladas, cada una con 256 unidades ocultas.
La salida de la primera capa recurrente (secuencia de estados ocultos) alimenta
la entrada de la segunda capa, permitiendo al modelo aprender representaciones
jerárquicas de mayor nivel de abstracción. Dependiendo de la configuración, las
capas pueden ser LSTM o GRU, manteniendo una interfaz idéntica para ambas variantes.

\subsection{Estrategia de Regularización}

Para prevenir el sobreajuste, el modelo implementa una estrategia de regularización
basada en dropout en dos puntos:

\begin{itemize}
    \item \textbf{Dropout post-embedding} ($p = 0.5$): aplicado inmediatamente después
    de la capa de embedding, antes de alimentar la secuencia a las capas recurrentes.
    Este dropout agresivo actúa como una forma de aumento de datos a nivel de
    representación.
    \item \textbf{Dropout entre capas RNN} ($p = 0.3$): aplicado entre las capas
    recurrentes apiladas mediante el parámetro \texttt{dropout} del módulo RNN de
    PyTorch. Este regulariza las representaciones intermedias entre capas.
\end{itemize}

\subsection{Capa de Clasificación}

La capa final es una transformación lineal que recibe el estado oculto de la última
capa recurrente en el último paso temporal, con dimensión $256 \to 1$. La salida
es un logit escalar que, al aplicarle la función sigmoide, produce la probabilidad
de que la reseña sea positiva. Esta capa contiene $256 + 1 = 257$ parámetros
(pesos más bias).

\subsection{Inicialización del Bias de la Compuerta de Olvido}

Para las variantes LSTM, el bias de la compuerta de olvido (\textit{forget gate})
se inicializa a 1.0 siguiendo la recomendación de \cite{jozefowicz2015}. Esta
inicialización favorece que la celda LSTM recuerde información desde las primeras
etapas del entrenamiento, evitando que el gradiente se desvanezca prematuramente
a través de la compuerta de olvido antes de que el modelo haya aprendido qué
información retener.


\subsection{Resumen de la Arquitectura}

\begin{table}[H]
\centering
\caption{Arquitectura completa del modelo \texttt{SentimentRNN}. Se muestran las
dimensiones de entrada y salida de cada capa, así como el número de parámetros
entrenables para las variantes LSTM y GRU.}
\label{tab:arquitectura}
\begin{tabular}{lllll}
\toprule
\textbf{Capa} & \textbf{Tipo} & \textbf{Entrada} & \textbf{Salida} & \textbf{Parámetros} \\
\midrule
Embedding       & \texttt{nn.Embedding}    & 25,000 & 128 & 3,200,000 \\
Dropout (embed) & \texttt{nn.Dropout(0.5)} & 128    & 128 & 0 \\
RNN Capa 1      & LSTM/GRU                 & 128    & 256 & LSTM: 394,240 / GRU: 295,680 \\
RNN Capa 2      & LSTM/GRU                 & 256    & 256 & LSTM: 525,312 / GRU: 393,984 \\
Lineal          & \texttt{nn.Linear}       & 256    & 1   & 257 \\
\bottomrule
\end{tabular}
\end{table}

\subsection{Implementación en Código}

A continuación se presentan los fragmentos clave de la implementación en
\texttt{model.py} que materializan la arquitectura descrita.

% --- Definición de capas del modelo (SentimentRNN.__init__) ---
El constructor de la clase \texttt{SentimentRNN} define las capas del modelo:
embedding, dropout, RNN apilada y capa lineal de salida. La selección entre
LSTM y GRU se realiza dinámicamente mediante el parámetro \texttt{rnn\_type}.

\begin{minted}{python}
class SentimentRNN(nn.Module):
    def __init__(
        self,
        vocab_size: int,
        embedding_dim: int,
        hidden_size: int,
        num_layers: int,
        rnn_type: str,
        rnn_dropout: float,
        embed_dropout: float,
    ) -> None:
        super().__init__()

        if rnn_type not in ("LSTM", "GRU"):
            raise ValueError(
                f"Invalid rnn_type '{rnn_type}'. "
                f"Supported options: 'LSTM', 'GRU'"
            )

        self.rnn_type = rnn_type
        self.hidden_size = hidden_size
        self.num_layers = num_layers

        # Embedding layer
        self.embedding = nn.Embedding(vocab_size, embedding_dim)

        # Dropout after embedding
        self.embed_dropout = nn.Dropout(p=embed_dropout)

        # Stacked RNN (LSTM or GRU)
        rnn_cls = nn.LSTM if rnn_type == "LSTM" else nn.GRU
        self.rnn = rnn_cls(
            input_size=embedding_dim,
            hidden_size=hidden_size,
            num_layers=num_layers,
            dropout=rnn_dropout,
            batch_first=True,
        )

        # Output linear layer
        self.fc = nn.Linear(hidden_size, 1)
\end{minted}

% --- Inicialización del bias de forget gate a 1.0 ---
La inicialización del bias de la compuerta de olvido a 1.0 se aplica únicamente
a las variantes LSTM. El bias del módulo LSTM de PyTorch se organiza como un
vector concatenado de cuatro compuertas, cada una de tamaño \texttt{hidden\_size}.
El segmento correspondiente a la compuerta de olvido ocupa las posiciones
$[n/4, n/2)$ del vector de bias.

\begin{minted}{python}
        # Initialize LSTM forget gate bias to 1.0 to encourage
        # remembering early in training (Jozefowicz et al., 2015)
        if rnn_type == "LSTM":
            for name, param in self.rnn.named_parameters():
                if "bias" in name:
                    # LSTM bias is organized as:
                    # [input_gate, forget_gate, cell_gate, output_gate]
                    # each of size hidden_size.
                    # Set forget gate bias (positions n//4 to n//2)
                    # to 1.0 for better gradient flow.
                    n = param.size(0)
                    param.data[n // 4 : n // 2].fill_(1.0)
\end{minted}

% --- Método forward: flujo de datos embedding → RNN → fc ---
El método \texttt{forward} implementa el flujo de datos completo: los índices de
tokens pasan por el embedding y dropout, luego por las capas recurrentes apiladas,
y finalmente el estado oculto de la última capa se proyecta a un logit escalar.

\begin{minted}{python}
    def forward(self, x: Tensor) -> Tensor:
        # Embedding: (batch, seq_len) -> (batch, seq_len, embed_dim)
        embedded = self.embedding(x)

        # Apply dropout after embedding
        embedded = self.embed_dropout(embedded)

        # RNN forward pass
        # output: (batch, seq_len, hidden_size)
        _output, hidden = self.rnn(embedded)

        # Extract hidden state from the final layer
        if self.rnn_type == "LSTM":
            # h_n: (num_layers, batch, hidden_size)
            final_hidden = hidden[0][-1]  # (batch, hidden_size)
        else:
            # h_n: (num_layers, batch, hidden_size)
            final_hidden = hidden[-1]  # (batch, hidden_size)

        # Linear: (batch, hidden_size) -> (batch, 1)
        logits = self.fc(final_hidden)

        return logits
\end{minted}
